\section{Modelling the Physical Problem}
This section introduces the mathematical framework and governing equations used to model the two-phase flow problems, laying the foundation for subsequent numerical simulations.
The the rationale behind averaging techniques for multiphase flow modelling is explained and the 
 macroscopic formulations used to represent complex two-phase dynamics introduced. The 'one-fluid' and 
 'two-fluid' models are detailed.


\subsection{Evaporation: Phenomenological Overview and Literature Review}
Evaporation processes are characterised based on whether they occur at saturation temperatures or not.
 Evaporation at saturation temperatures, known as vaporisation, is a much faster 
 process in comparison to evaporation at normal~\cite{wickertSimulation}. 
 In cases where vaporisation occurs 
 along a heated wall, most CFD models distribute the energy 
 transfer by separating the wall heat flux into two components: 
 latent heat absorption associated with bubble vaporisation and 
 sensible heat transferred to the liquid through conduction, 
 convection, and quenching processes. The partitioning depends on 
 the area fractions covered by liquid and vapour phases.
 Empirical results or models of subprocesses have 
 established correlations between the heat flux and wall superheat.

The sub-processes that are primarily taken into account in such determinations are:
\begin{itemize}
    \item number density of nucleation sites,
    \item diameter of bubble at time of detachment from the surface,
    \item and the frequency of bubble detachment.
\end{itemize}
Nucleation affects the number and distribution of bubbles in the fluid. Walls, particularly so in case they're heated, are more susceptible to be nucleation sites for bubble growth
 as the activation energy for bubble growth is much lower. This nucleation can be modelled through the domain as homogenous or heterogenous - where nucleation sites are modelled using a 
 presumed PDF function or be determined by empirical correlations.~\cite{liaoComputationalModellingFlash2017}

After a stable bubble nucleus has been formed, the growth of the bubble becomes pertinent in the evaluation
 of the diameter at departure and the interphase energy and mass transfer. In the case of an ideal
 spherical bubble undergoing expansion in a uniformly superheated liquid, the rate of growth
 is determined by the surface tension, liquid inertia, and the difference in pressures inside
 and outside the bubble. Surface tension dominates and restricts the initial bubble growth.
 This initial rapid expansion of the bubble results in the formation of a thin liquid layer separating
 the vapour from the wall surface. This microlayer can be a significant contributor to the 
 bubble growth and heat transfer process.
 As bubble radius increases, the liquid layer surrounding the bubble 
 cools due to evaporation, resulting in a decrease in vapour pressure. 
 The bubble growth dynamics are initially governed by thermal diffusion 
 and inertia. Over time, inertial effects become less important as 
 the bubble expands more slowly. Further growth is determined 
 primarily by interphase energy exchange.~\cite{sherFlashboilingAtomization2008}

\subsubsection{Effects of low pressure and gravity conditions}
The absence of significant gravitational forces alters the dynamics of vaporisation - affecting natural convection, buoyancy
 driven phenomenon, and bubble dynamics. Under terrestrial gravity, convection currents develop under thermal
 gradients, thus bringing cooler liquid to the heated surface as the vapour bubble rises due
 to buoyoancy. This density driven natural convection is affected in low gravity conditions, maintaining
 the thermal gradient. Bulk fluid movement is affected and heat and mass transfer processes shift
 from being convection dominated to being diffusion dominated.~\cite{cochranEffectsFluidProperties, leeExperimentalComputationalInvestigation2022}

Under low gravity conditions, vapour bubbles tend to remain near the heat source, growing
 significantly larger and coalescing with smaller bubbles. Prolonged bubble presence on the
 surface results in drying out of the surface and unsteady rise in surface temperatures.\cite{qiuSingleBubbleDynamicsPool2002}~\COMMENT{Mention some numbers}
 In conditions where the dominance of surface tension forces becomes more pronounced,
 the surface tension gradients can induce fluid circulation. This is known as Marangoni
 convection and is observed to affect bubble growth and heat transfer.~\cite{xuNumericalStudyMarangoni2014}

Reduced system pressure impacts the overall intensity of heat transfer. The onset of nucleate boiling also
 requires higher superheat values in low pressure conditions.
 Decrease in pressure leads to a lower number of active nucleation sites and increases
 microlayer thickness found at the bubble-wall interface. The local heat flux density from
 microlayer evaporation is impacted and the microlayer evaporation area and depletion time is
 increased. The system is also susceptible to the formation of a continuous vapour film at the 
 surface due to coalescing bubbles unable to escape the surface, reducing the overall heat transfer.~\cite{serdyukovInfluencePressureLocal2023}
 
The affected heat transfer and bubble growth phenomenon further result in lower critical heat fluxes
 and heat transfer efficiency.~\cite{zhangExperimentalTheoreticalStudy2002, PDFMethodAssessing}

\subsubsection{Vaporisation in a superheated liquid}
Flash evaporation is the intense vaporisation triggered through depressurisation 
 of superheated fluids below the saturation pressure.
 This process induces a violent and transient heat
 and mass transfer, resulting in the explosive formation and growth
 of vapour bubbles throughout the liquid volume. The phenomenon is significant 
 in desalination on ships, nuclear reactor safety, and venting of cryogenic
 propellant tanks, and the harvesting of low-temperature geothermal energy.~\cite{sauryExperimentalStudyFlash2002, wuExperimentalStudyFlash2019, liaoReviewNumericalModelling2021, taniNUMERICALSIMULATIONSINGLE2025, mutairExperimentalInvestigationCharacteristics2010, vahajiEfficiencyTwophaseNozzle2014}.

Early experiments that systematically studied flash evaporation in liquid pools and sprays in superheated 
 water jets and early empirical correlations to determine the amount of generated 
 vapour were established~\cite{miyatakeExperimentalStudySpray1981, miyatakeEffectLiquidTemperature1981, gopalakrishnaExperimentalStudyFlash1987, petersonInvestigationsLiquidflashingEvaporation1984}. 
 For flash evaporation in static pools, the amount of liquid height increased the duration 
 of the phenomenon and the rate of depressurisation increases the intensity.
 The final evaporated mass also increases proportional to the amount of superheat.
 While greater thermal inertia at larger depths slows the flashing process, the amount of total 
 mass evaporated does not increase proportionally.~\cite{sauryFlashEvaporationWater2005,sauryExperimentalStudyFlash2002}
 It has been observed during flashing, that a liquid-vapour interface is generated 
 on the free surface which then progresses into the superheated liquid~\cite{HahneBarthau2000, reinkeExplosiveVaporizationSuperheated2001}.
 High-speed imaging and temperature measurements consistently show that
 flash evaporation proceeds through three distinct stages: bubble generation,
 boiling growth, and a boiling reduction stage~\cite{wuExperimentalStudyFlash2019a}.
 For the final evaporated mass, under moderate superheat conditions $\Delta T \leq 35$\,K,
 a strong linear scaling is observed~\cite{liaoFlashingEvaporationDifferent2013,liaoReviewNumericalModelling2021}.
 Recent experiments studying stationary volatile drop evaporation in microgravity
 report a $56\%$ reduction in evaporation rate as compared to normal gravity conditions~\cite{kumarSessileVolatileDrop2020}.
 ~\COMMENT{need to add conclusion to this that ties into choice of phase change model later} 
%The above mentioned studies have proposed various correlations to quantify the flash evaporation phenomenon.
% The bulk temperature response during flash evaporation can be quantified by the 
% non-equilibrium function (NEF),
% \begin{equation}\label{eq:flashing-nef}
%   NEF = \frac{T(t) - T_{e}}{T_{e} - T_{0}}
% \end{equation}
% Correlations between the evaporated mass at time $t$ and the final evaporated mass 
% have been proposed.~\COMMENT{list em maybe}

\subsection{Modelling Multiphase Flows}
The computational tribulations of directly solving local instantaneous governing equations throughout 
 the entire domain necessitates dimensional reduction through macroscopic formulations. 
 For two-phase flows (used interchangeably with multiphase flows), these formulations rely 
 on averaging local flow field fluctuations while preserving continuum mechanics principles. 

In the following, we use the subscript $i \in \{l, v\}$ to denote variables associate with the 
 liquid (l) and vapour (v) phase respectively. Figure~\COMMENT{FIG} displays a region $\Omega$ where two phases $l,v$, or equivalently $p,s$ (primary and secondary), occupy
 regions $\Omega_{l}(t)$ and $\Omega_{v}(t)$ such that $\Omega = \Omega_{l}(t) \cup \Omega_{v}(t)$ is constant in time.    
 The boundaries of the domains - $\partial \Omega_{l}$, $\partial \Omega_{v}$, and $\partial \Omega$ - together define the interface that separates the two phases as
 \begin{equation}
  \mathcal{S}(t) \;=\;\Bigl(\partial\Omega_l\cap\partial\Omega_v\Bigr)\setminus\partial\Omega\,.
  \label{eq:interface}
\end{equation}~\COMMENT{Equation seems a bit pretentious or? also idk if is a good choice to make here for the notation}
 Eulerian averaging is related to the Eulerian description of the flow - the local instantaneous
 variables are averaged with the time and space variables being kept constant.
 Volume averaging admits a general 
 representative volume element $V$ bounded by $\partial V$, and 
 time averaging uses a time interval at
 every spatial reference point $(\vetor{x}, t)$ in the region.
 The averaging volume $V$ is partitioned into $V_i$ for phase $i\in\{l,v\}$ with phase volume fraction
\begin{equation}
  \alpha_i(\vetor{x},t) \;=\; \frac{V_i}{V}\,,\qquad
  \alpha_l + \alpha_v = 1\,.
  \label{eq:volfrac}
\end{equation}

All averaging procedures generally utilise a phase indicator function that 
 marks the presence of the phase in the averaging domain. Denoted as $\gamma_{i}$, 
 the phase indicator function can be used to define the average of the general scalar or vector quantity $\psi$ over the averaging 
 volume $V$,
\begin{equation}
  \overline{\psi_i}^V
  \;=\;
  \frac{1}{V}\int_V \gamma_i\,\psi_i(\vetor{x},t)\,\mathrm{d}V\,,
  \label{eq:volavg}
\end{equation}
and the phase-intrinsic average then is
\begin{equation}
  \overline{\psi_i}^i \;=\;\frac{1}{V_i}\int_{V_i}\psi(\vetor{x},t)\,\mathrm{d}V
  \;=\;
  \frac{\overline{\psi_i}^V}{\alpha_i}\,.
  \label{eq:intravg}
\end{equation}

Substituting instantaneous variables with their averaged 
 equivalents acts as a low-pass filter, smoothing out discontinuities 
 at the phase interface. This process eliminates direct knowledge of 
 interface topology, thus necessitating closure models and 
 explicit interfacial transfer terms. Notably, while Star-CCM+ 
 employs Euler-averaged equations for its multiphase models, 
 the software documentation does not specify which averaging 
 approach - time, volume, or a combination - is used.
A thorough treatise on the averaging procedures can be found~\COMMENT{here and here}. The development
of the two-fluid model and one-fluid model for two-phase flows is detailed~\COMMENT{here and here} respectively. 
 ~\COMMENT{language}~\cite{ishiiThermoFluidDynamicsTwoPhase2011,wornerCompactIntroductionNumerical, prosperettiCOMPUTATIONALMETHODSMULTIPHASE, brennenFundamentalsMultiphaseFlow2006}~\citelater{drew, star-ccm+}


Formulations of the macroscopic field equations for multiphase flows can either treat each phase
 separately or consider the mixture as a whole. When the dynamics
 of the two phases are closely coupled, the two sets of governing equations
 can be summed up together which then along with the relevant interfacial transfer 
 terms describe the behaviour
 of the mixture. In the two-fluid formulation, the set of the averaged governing 
 equations is solved for each phase. The phases can evolve independently and are
 coupled through interfacial transfer terms.
 In both formulations, the transfer terms coupling the phases must also 
 account for the relevant jump conditions.~\cite{ishiiThermoFluidDynamicsTwoPhase2011,wornerCompactIntroductionNumerical,tryggvasonDirectNumericalSimulations2011, yadigarogluNatureMultiphaseFlows2018}

Often a 'one-fluid' model is utilised to model separate flows due to its ease of implementation
 and computational efficiency. The phases are treated as one fluid with variable 
 material properties that have an abrupt change at the interface. Interface capturing 
 methods such as the volume-of-fluid (VOF) employ a marker function that distinguishes
 the phases and is advected with the evolution of the flow. This marker function can be used
 to reconstruct the interface. The governing equations are analogous to the 
 mixture formulation of two-phase flows with no relative velocity between the phases.~\cite{gadaLargeScaleInterface2017,tryggvasonDirectNumericalSimulations2011,prosperettiCOMPUTATIONALMETHODSMULTIPHASE}~\citelater{star-ccm+}

The following discussion presumes that the averaging procedure has been performed and discusses the governing equations
in their averaged macroscopic form as solved in Star-CCM+. The quantities are assumed to be 
the phase intrinsic averages. For both the formulations, we assume the validity of the continuum 
hypothesis, that the interfaces are infinitely thin, and neglect intermolecular forces.


\subsubsection{The One-Fluid Formulation}\label{sec:one-fluid}
In order to account for the varying material properties when representing
 two-phase flow as a single fluid, the mixture properties are defined in the 
 volume element $V$. 
 The mixture density $\rho$, viscosity $\mu$, and thermal conductivity $k$ 
 are defined as
 \begin{subequations}
  \begin{align}
    \rho &\;=\;\textstyle\sum\nolimits_{i} \alpha_{i} \rho_{i}\,,\\
    \mu &\;=\; \textstyle\sum\nolimits_{i} \alpha_{i} \mu_{i}\,\\
    k &\;=\; \textstyle\sum\nolimits_{i} \alpha_{i} k_{i}\,.
  \end{align}
\end{subequations}
 Here, the properties $\rho_{i}$, $\mu_{i}$, and $k_{i}$ are phase intrinsic averages.
 Using the defined mixture density $\rho$, the velocity, specific heat, and enthalpy can 
 be also defined as mass averaged mixture values.
 \begin{subequations}\label{eq_vof_defs}
\begin{align}
  \vetor{v} &\;=\; \textstyle\sum\nolimits_{i} \alpha_{i} \rho_{i} \vetor{v}_{i}/ \rho,\\
  C_{p} &\;=\; \textstyle\sum\nolimits_{i} \alpha_{i} C_{p,i} \rho_{i}/\rho,\\
  H &\;=\; \textstyle\sum\nolimits_{i} \alpha_{i} \rho_{i} \mathrm{H}_{i}/\rho + |\vetor{v}|^{2}/2\,,\text{where $\vetor{v}$ is the mixture velocity}. 
\end{align}
\end{subequations}
Sometimes a relative velocity between the phases is considered, 
  the diffusion velocity $\vetor{v}_{d,i}$ of phase $i$ can be defined as
\begin{equation}
  \vetor{v}_{d,i} \;=\; \vetor{v}_{i} - \vetor{v}\,.
\end{equation}
The stress tensor $\tens{T}$ can be extended to include both the molecular and turbulent stresses. 
 The various forces are accounted through the pressure $p$, gravitational acceleration $\vetor{g}$,
 and the body force vector $\vetor{f_{b}}$. $\tens{I}$ is the second order identity tensor.
 The total energy of the system is given as 
 $E$ and the heat flux vector is denoted by $\vetor{\dot{q}''}$.  
Given mass, momentum, and energy sources for the phase $i$ as $S_{i}^{\alpha}$, $\vetor{S}_{i}^{\alpha}$, and
 $S_{i}^{E}$ respectively, the governing equations for the one-fluid formulation are given 
 as follows: 
\begin{subequations}\label{eq_vof_mass}
\begin{align}
\frac{\partial}{\partial t}\int\limits_{V} \rho\,\mathrm{d}V +
\int\limits_{\partial V} \rho \vetor{v} \cdot \mathrm{d}\vetor{a}
&\;=\; \int\limits_{V} \sum_{i}\mathrm{S}^{\alpha}_{i} \rho_{i}\,\mathrm{d}V\,,\\
%
\frac{\partial}{\partial t}\int\limits_{V} \rho \vetor{v}\,\mathrm{d}V +
\int\limits_{\partial V} \rho\vetor{v} \otimes \vetor{v} \cdot \mathrm{d}\vetor{a}
&\;=\; \int\limits_{\partial V} \left(-p \tens{I} + \tens{T}\right)\cdot\mathrm{d}\vetor{a} + \int\limits_{V} \rho \vetor{g}\,\mathrm{d}V
 + \int\limits_{V} \vetor{f}_{b}\,\mathrm{d}V \nonumber\\
&\qquad\quad{} - \sum_{i}\int\limits_{\partial V} \alpha_{i} \rho_{i} \vetor{v}_{d,i} \otimes \vetor{v}_{d,i} \cdot\mathrm{d}\vetor{a}
 + \int\limits_{V} \sum_{i}\vetor{S}^{\alpha}_{i}\,\mathrm{d}V\,\\
%
\frac{\partial}{\partial t}\int\limits_{V}\rho \mathrm{E}\,\mathrm{d}V +
\int\limits_{\partial V} \left[\rho H \vetor{v} + \sum_{i}\alpha_{i} \rho_{i} \mathrm{H}_{i} \vetor{v}_{d,i}\right] \cdot \mathrm{d}\vetor{a}
&\;=\; - \int\limits_{\partial V} \vetor{\dot{q}''}  \cdot \mathrm{d}\vetor{a} + \int\limits_{\partial V} \tens{T}\cdot\vetor{v}\,\mathrm{d}\vetor{a}
 + \int\limits_{V}\vetor{f}_{b}\cdot\vetor{v}\,\mathrm{d}V\nonumber\\
&\qquad{} + \int\limits_{V} \sum_{i} S_{i}^{E}\,\mathrm{d}V
\end{align}
\end{subequations}

The computational efficiency and ease of implementation of the one-fluid formulation is limited by increased coupling of the 
 phases and complex interfacial dynamics - as the averaging volume must be 
 able to capture the various scales involved in the interphase interactions. As jump conditions are not explicitly enforced, interfacial terms must be added 
 to the governing equations at the interface which brings in numerical complications. Many existing one-fluid formulations 
 in the literature are analytically inconsistent under vaporization. Under conditions of phase change, 
 the conservative and non-conservative forms of the governing equations for various formulations are found to not be equivalent 
 and that the momentum and energy jump conditions are not recovered.~\cite{trujilloReexaminingOnefluidFormulation2021}. Phase
 change is generally modelled using source terms in the volume fraction equation and the energy equation~\cite{hardtEvaporationModelInterfacial2008,tryggvasonDirectNumericalSimulations2011}.

\subsubsection{The Two-Fluid Formulation}
Considering the volume element $V$ as described in the previous section, 
 the volume fractions of the two phases $\alpha_{i}\,i \in \{l,v\}$ must satisfy the continuity constraint 
 \begin{equation}\label{eq:vol-frac-cont}
  \sum_{i}\alpha_{i}\;=\;1\,.
 \end{equation}
 The volume of the phase $i$ is given as 
 \begin{equation}
  V_{i} = \int\limits_{V} \alpha_{i}\,\mathrm{d}V\,,
 \end{equation}
 The ratio $V_{i}/V$ of the effective volume of the phase $i$ to the total volume is defined as 
 the void fraction of the phase. 
 
The averaged macroscopic equation for the conservation of mass is given in~\eqref{eq:emp-mass}.
 The density and velocity of the phase $i$ is denoted by $\rho_{i}$ and $\vetor{v}_{i}$.
 The interfacial mass transfer unit volume between the two phases is denoted
 by $\Gamma$ and the surface area vector of the 
 volume element is denoted by $\vetor{a}$. With an additional phase mass source term $\mathcal{S}_{i}^{\alpha}$,
 the conservation of mass for a generic phase $i$ is 
 \begin{equation}\label{eq:emp-mass}
\frac{\partial}{\partial t}
\int\limits_{V} \alpha_{i} \rho_{i} \,\mathrm{d}V
+ \int\limits_{\partial V} \alpha_{i} \rho_{i} \vetor{v}_{i} \cdot \mathrm{d}\vetor{a}\;=\;
\int\limits_{V} \Gamma_{i} \,\mathrm{d}V
+ \int\limits_{V} \mathrm{S}^{\alpha}_{i} \,\mathrm{d}V\,.
\end{equation}
The interfacial mass transfer term $\Gamma_{i}$
 represents the rate of mass transfer from the phase $i$. The interaction area density $a_{lv}$ specifies 
 the interfacial area available in the volume element over which the mass transfer occurs.
\begin{equation}
 \Gamma_{i}\;=\;a_{lv}\;\dot{m}_{i}\;=\; \rho_{i} \vetor{n}_{i} \cdot (\vetor{v}_{i} - \vetor{v}_{\mathcal{S}})\,.
\end{equation}
The term $\rho_{i} \vetor{n}_{i} \cdot (\vetor{v}_{i} - \vetor{v}_{\mathcal{S}})$ consists of the local instantaneous 
 variables instead of the averaged quantities. The interface normal vector and interface 
 velocity are given as $\vetor{n}_{i}$ and $\vetor{v}_{\mathcal{S}}$ respectively. 


The governing equation for the conservation of momentum is given in~\eqref{eq:emp-mom}.
 The equation describes the change in momentum of the phase $i$ due to the effects of the
 fundamental forces in the system along with the effects of any internal forces or potentials $(\vetor{F_{int}})_{i}$
 or any phase momentum source terms $\vetor{S}_{i}^{\alpha}$. The pressure and gravitational acceleration is given as $p$ and $\vetor{g}$ respectively.
 The stress tensor $\tens{T}_{i}$ comprises of the molecular and turbulent stress tensors. Turbulence
 in addressed in~\ref{sec:turbs}.
\begin{equation}\label{eq:emp-mom}
\begin{split}
\frac{\partial}{\partial t}
\int\limits_{V} \alpha_{i} \rho_{i} \vetor{v}_{i} \,\mathrm{d}V
+ \int\limits_{\partial V} \alpha_{i} \rho_{i} \vetor{v}_{i} \otimes \vetor{v}_{i} \cdot\,\mathrm{d}\vetor{a}
&\;=\; 
\int\limits_{V} \left(-\alpha_{i} \nabla p 
+ \alpha_{i} \rho_{i} \vetor{g}
+ (\vetor{F_{int}})_{i} \right) \,\mathrm{d}V
+ \int\limits_{\partial V} \left[ \alpha_{i} \left( \tens{T}_{i}\right) \right] \cdot \,\mathrm{d}\vetor{a} \\
&\qquad {} + \int\limits_{V} \vetor{S}^{\alpha}_{i} \,\mathrm{d}V
+ \int\limits_{V}  \vetor{M}_{i}  \,\mathrm{d}V\,.
\end{split}
\end{equation}
The momentum transfer term $\vetor{M}_{i}$ accounts for the momentum transfer from 
 the phase $i$ to the interface. This term also accounts for the momentum transfer
 due to interfacial mass transfer and is further detailed in Section~\ref{sec:closures}. It is only non-zero in case 
 of a relative motion between the phases and is subject to the constraint
 \begin{equation}
  \vetor{M}_{i} \;=\; a_{lv}\left(\left(-p_{i} + \tens{T}_{i}\right)\cdot\vetor{n}_{i} - \dot{m}_{i}\vetor{v}_{i}\right)\,.
 \end{equation}
In case of constant surface tension, the dynamic condition that the effects of the hydrodynamic forces are 
 in equilibrium with the forces due to surface tension. Denoting the part of the interface $\mathcal{S}$ in the 
 volume element $V$ as $\mathcal{S} \cap V$, the dynamic jump condition can be written as
 \begin{equation}
  \vetor{M}_{l} + \vetor{M}_{v}\;=\; \frac{1}{V}\int\limits_{\mathcal{S} \cap V} 
    \vetor{f_{\sigma}}\,\mathrm{d}V\,.
 \end{equation}

%The rate of change of momentum due to the 
% interfacial mass transfer can be given using the velocities of the phases $l$ and $v$ and the rate 
% of mass transfer between the two phases,
%\begin{equation}
%\vetor{\Gamma} = (m_{vl}\vetor{v}_{l} - m_{lv}\vetor{v}_{v}) a_{lv}\,.
%\end{equation}

The averaged macroscopic formulation of the conservation of energy in the two-fluid formulation is 
 shown in~\eqref{eq:emp-en}. The equation accounts for the change of the total energy $E_{i}$ in the 
 volume $V$ due to the work done by line, surface, and volume (body) forces on the volume element,
 any heat transfer (in this specific case, only conduction and convection are considered),
 and due to the interfacial mass transfer in the volume element. 
  
\begin{align}\label{eq:emp-en}
\frac{\partial}{\partial t}
\int\limits_{V} \alpha_{i} \rho_{i} E_{i} \,\mathrm{d}V
+ \int\limits_{\partial V} \alpha_{i} \rho_{i} H_{i}&
\vetor{v}_{i} \cdot \mathrm{d}\vetor{a}
\;=\;
\int\limits_{\partial V} \alpha_{i} (k_{\mathrm{eff}})_{i} \nabla T_{i} \cdot \mathrm{d}\vetor{a}
+ \int\limits_{\partial V} \tens{T}_{i} \vetor{v}_{i} \cdot \mathrm{d}\vetor{a}\nonumber \\
& + \int\limits_{V} \vetor{f_{b}}_{i} \cdot \vetor{v}_{i} \,\mathrm{d}V
+ \int\limits_{V} \mathrm{S}_{u,i} \,\mathrm{d}V + \Lambda.
\end{align} 

%[ 
%    \underbrace{
%        h_{i}^{\text{REF}} 
%+ \int\limits_{T_{i}^{\text{REF}}}^{T_{i}} C_{p,i}(T') \,\mathrm{d}T' 
%+ \frac{1}{2} \vetor{v}_{i} \cdot \vetor{v}_{i}
%}_{\text{total enthalpy, $H_{i}$}}
%]
Having each transfer process of each phase balanced separately, 
 the two-fluid formulation is much more complex and requires closure models
 for each interaction. However, the two-fluid formulation is a much more
 complete and general description of two-phase flow - accounting for the dynamic and no-equilibrium 
 interactions between the phases. The mixture or the two-fluid model do not retain information about the interface topology.
 Given a certain void fraction of a phase the two models are
 unable to comment on the nature of the flow regime.~\cite{ishiiThermoFluidDynamicsTwoPhase2011}.

\subsection{Modelling Turbulence}\label{sec:turbs}
Flows are generally turbulent. Turbulence is characterised by irregular fluid motion dominated by
 eddies and vortices across a wide range of spatial and temporal scales. The largest eddies
 extract kinetic energy from the mean flow, which is then, through a cascade, transferred
 to the eddies at
 the smallest scale where the energy is dissipated due to viscosity, raising the 
 internal energy of the fluid~\cite{gilmanDevelopmentGeneralPurpose2014, davidsonIntroductionTurbulenceModels}.

RANS models result from decomposing flow variables into
 mean and fluctuating components, $\displaystyle \phi = \overline{\phi} + \phi'$, applying Reynolds averaging,
 which filters out turbulent fluctuations below the
 averaging scale~\cite{besnardTurbulenceMultiphaseFlow1988a, hirschNumericalComputationInternal2007}.
 This operation yields additional Reynolds stress terms representing
 momentum transport by turbulence, which require closure.
 In order to draw closure relations, eddy-viscosity models of turbulence assume that:
 \begin{itemize}
    \item the turbulent scalar flux vector aligns with the mean scalar gradient.
    \item the Reynolds stress tensor is proportional to the mean velocity strain rate through an eddy viscosity.
 \end{itemize}
 While these assumptions simplify modelling, they tend to be, more often than not, invalid
 as they do not fully capture turbulence physics, especiialy in anisotropic 
 or non-equilibrium flow~\cite{popeTurbulentFlows2015}. This limitation can be significant
 in regions with complex strain rates or strong secondary flows.

Although turbulence can be modelled in various ways to various levels of resolutions,
 turbulence in this study has been modelled using the Shear Stress Transport (SST) k-$\omega$ model.
 The SST k-$\omega$ model is a two-equation eddy-viscosity turbulence model
 based on the Reynolds-Averaged Navier-Stokes (RANS) equations.
 The model combines the advantages of the $k$-$\omega$ model near walls - which
 accurately resolves boundary layer behaviour - and the $k$-$\varepsilon$ model in
 the free stream, enhancing robustness and accuracy in separated flows~\cite{wilcoxTurbulenceModelingCFD2006}.
 The SST approach uses a blending function to smoothly transition between these models,
 which is especially effective in flows with strong adverse pressure gradients and separation,
 such as multiphase cavity flows~\cite{Menter1994}. Another blending function limits the eddy viscosity 
 in wake regions and a cross-diffusion term ensures consistency between the $k$-$\omega$ and $k$-$\epsilon$
 models.

Since the form of the RANS governing equations remains similar to the governing equations described in the
 previous section, they are not explicitly written here but can be found in the Star-CCM+ documentation~\citelater{starccm}.
 The additional equations accounting for the transport of the turbulent kinetic energy and dissipation rate 
 is given in~\ref{eq:turbs-kw}
 \begin{subequations}\label{eq:turbs-kw}
 \begin{align}
    \frac{\partial \rho k}{\partial t}  + \nabla (\rho k \vetor{\overline{v}}) &= \nabla \left[(\mu + \sigma_{k} \mu_{t})\nabla k\right] + P_k
        - \rho \beta^{*}f_{\beta^{*}} (\omega k - \omega_0 k_0) + S_k,\\
    \frac{\partial \rho \omega}{\partial t} \nabla (\rho \omega \vetor{\overline{v}}) & = \nabla\left[(\mu + \sigma_{\omega}\mu_t) \nabla \omega\right] + P_{\omega}
        - \rho \beta f_{\beta}(\omega^{2} - \omega_{0}^{2}) + S_{\omega},
 \end{align} 
 \end{subequations}
 The turbulent viscosity, $\mu_t$, is calculated as $\mu_{t} = \rho k T{t}$, where $T_{t}$ is the turbulent time scale. 
 where $\sigma_{k}$ and $\sigma_{\omega}$ are the Prandtl numbers for $k$ and $\omega$ respectively. 


\subsection{Modelling Subproblem 1}
Not considering the mixing effects that might occur at later times, the infiltrating
 water into the evaporation chamber is a separated flow problem due to the sharp interface
 separating the two phases - water and air.
 
The objective - we want to understand the development of the flow field, 
 understand the amount of sloshing that can be expected in the system, and the amount of
 mixing that happens during the injection process and what the result flow system might look like 
 before the evaporation step. 

The fundamental forces active in the system are the inertial, gravitational, viscous,
 pressure, and surface tension forces. The governing equations are analogous to the 
 governing equations of the homogenous two-fluid flow model\footnote{two-fluid formulation with the phase velocities assumed to be constant}.
 The flow field is described using the primary variables for each phase - $\rho_{i}, \vec{u}_{i}, E_{i}$.

\begin{itemize}
   \item what we're looking for and why we're interested in this
   \item phases involved, forces active and shit
   \item assumptions we make
   \item variables we care about
   \item governing eqs (?)
   \item 3D vs 2D
\end{itemize}

\subsection{Modelling Subproblem 2}
\begin{itemize}
  \item from the top 
\end{itemize}